\documentclass[11pt]{article}
\usepackage{amsmath,amsthm}
\usepackage{named}
\usepackage{changepage}

\newtheorem{proposition}{Proposition}
\newtheorem{lemma}{Lemma}

\def\o{\overline}
\def\bs{\backslash}
\def\std{\mathit{Std}}
\def\head{\mathit{Head}}
\def\body{\mathit{Body}}
% \def\p2f{\mathit{p2f}}
\def\ph{\mathit{PH}}
\def\asm{\mathit{Asm}}
\def\ug{\mathit{UG}}
\def\HT{\mathit{HT}}
\def\fnt{\mathit{finite}}
\def\gringo{{\sc gringo}}
\def\clingo{{\sc clingo}}
\def\mg{mini-\gringo}
\def\ar{\leftarrow}
\def\rar{\rightarrow}
\def\uc{\widehat\forall}
\def\lrar{\leftrightarrow}
\def\beq{\begin{equation}}
\def\eeq#1{\label{#1}\end{equation}}
\def\ba{\begin{array}}
\def\ea{\end{array}}
\def\kbugsd{\emph{KB, UG, SD}}
\def\bfx{{\bf X}}
\def\sd{\mathit{SD}}
\newcommand{\co}{\ensuremath{\mathit{count}}}
\newcommand{\no}{\ensuremath{\mathit{not}}}

\hyphenation{lif-schitz}

\title{Representing Rules with Counting by Formulas\\
  (Incomplete Draft)}
% \author{Vladimir Lifschitz\\ University of Texas}

\begin{document}
% \date{}
\maketitle

\begin{adjustwidth}{-1cm}{-1cm}

   This note is an attempt to extend the definition of
   the ``natural translation''~$\nu$ \cite{lif21,lif23,fan26b,fan26}
 to programs with counting.
  
  \section{Syntax of rules}\label{sec:nnf}
  
  In this section,
  the syntax of rules over a \mg\ signature %\cite[Section~2.1]{fan26b}
  is   generalized to programs with counting.
    Mini-\gringo\ signatures, terms, atoms,
  literals, comparisons, \mg\ dialects and definite/indefinite function
  symbols as in the Dialects draft \cite[Sections~2 and~3]{fan26b}.  By~$D$ we
  denote a dialect of \mg\ that includes
  \begin{itemize}
  \item no basic symbols other that numerals and symbolic constants, and
    \item no indefinite function symbols
  other than the interval symbol $..$ characterized by the condition
  $$\widehat{..}(m,n)=\{k: m\leq k\leq n\}.$$
\end{itemize}
(These simplifying assumptions are satisfied for the N-Queens programs
\cite[(1)--(5) and (6)--(10)]{lif26}).
A term is \emph{interval-free} if it does not
contain the interval symbol.

A term of the form $\mathit{blank}(t_1,\dots,t_k)$, where $k\geq 2$,
can be abbreviated as $(t_1,\dots,t_k)$.

  An \emph{aggregate element} is a pair
$\bfx:{\bf L}$,
where $\bfx$ is a non-empty tuple of distinct variables,
and $\bf L$ is a conjunction of
literals and comparisons such that every member of~$\bfx$ occurs in~$\bf L$
\cite[Section~4]{lif22a}.  An \emph{aggregate atom} is an expression of
one of the forms
$$t_1\leq \co \{E\},\ \co \{E\}\leq t_2,\ t_1\leq \co \{E\}\leq t_2,$$
where~$E$ is an aggregate element, and~$t_1$, $t_2$ are interval-free terms.

A \emph{rule} is an expression of the form
\beq
H \ar B_1\land\cdots\land B_n
\eeq{rule}
($n\geq 0$), where the head~$H$ is either an atom or empty, and
each member~$B_i$ of the body is either a literal, or a comparison,
or an aggregate atom possibly preceded by \no.  (Since negations in front of
aggregate atoms are allowed, this is more general than in some of the
earlier work \cite{lif22a,fan24}.)

A variable that occurs in a rule~$R$ is \emph{local} in~$R$ if each of
its occurrences is within an aggregate element, and \emph{global} otherwise.
A rule is \emph{pure} if, for every aggregate element $\bfx:{\bf L}$ in its
body, all variables in the tuple~$\bfx$ are local \cite[Section~4]{lif22a}.

A \emph{program} is a finite set of pure rules.

An \emph{equation} is a comparison of the form $t_1=t_2$.
A rule is \emph{in numeric normal form
  (NNF)} if
\begin{itemize}
\item[(a)]
  symbolic constants do not occur in it in the scope
  of arithmetic function symbols other than the interval symbol, and
\item[(b)]
  every occurrence of a term of the form $t_1\,..\,t_2$ in it
is the right-hand side of an equation.
\end{itemize}
(So symbolic constants are allowed as endponts of an interval; this is
different from what the Dialects draft requires \cite[Section~3]{fan26b}.)
% All rules (1)--(10) in the N-Queens programs are in numeric normal form
% in the sense of this definition.

%\paragraph{Semi-normal rules}
A rule is called {\em semi-normal} if it satisfies condition (b).  Any semi-normal rule
can be transformed into an NNF rule as
follows:
\begin{itemize}
\item make the list $c_1,\dots,c_k$ of 
the symbolic constants that have at least one occurrence in the scope of an arithmetic operation other than interval $..$;
\item  substitute for them fresh variables $X_1,\dots,X_k$, and append $X_1=c_1 \wedge\dots\wedge X_k=c_k$ to the body.  
\end{itemize}

\section{The target language}

By~$\sigma^\#$ we denote the many-sorted signature that consists of
\begin{itemize}
\item [(i)] the sort \emph{general}, the sort \emph{set}, and the
  subsort \emph{integer} of \emph{general},
\item [(ii)] all numerals of~$D$ as object constants of the
  sort \emph{integer};
\item [(iii)] all symbolic constants of~$D$ as object
   constants of the sort \emph{general};
 \item [(iv)] expressions $c\bs k$, where~$c$ is a symbolic constant of~$D$
   and~$k$ is a positive integer, as $k$-ary function constants, with
   both arguments and value of the sort \emph{general};
\item [(v)] all definite arithmetic
  function symbols of~$D$ as function constants with both
  arguments and value of the sort \emph{integer};
\item[(vi)] the symbol $\co$ as the unary function symbol with the argument
  of the sort \emph{set} and the value of the sort \emph{general};
\item [(vii)] expressions $c/k$, where~$c$ is a symbolic constant of~$D$
  and~$k$ is a nonnegative integer, as $k$-ary predicate constants
  with arguments of the sort \emph{general};
\item [(viii)] the symbol $\leq$ as a binary predicate constant with both
  arguments of the sort \emph{general};
\item [(ix)] the symbol $\in$ as a binary predicate constant with the
  arguments of the sorts \emph{general} and \emph{set}.
\end{itemize}
(In comparison with $\sigma(D)$ defined in the Dialects draft
\cite[Section~4.1]{fan26b}, three things here
are new: the sort \emph{set}, the function symbol $\co$, and the
predicate symbol~$\in$.  On the other hand, the
symbol~$\mathit{graph}_{..}$ is not included.)

A term of the form $(\mathit{blank}\bs k)(t_1,\dots,t_k)$, where $k\geq 2$,
can be abbreviated as $(t_1,\dots,t_k)$.


Inequalities involving order relations other that~$\leq$, as well as chains
of inequalities, are treated as abbreviations. We identify general variables
with variables occurring in rules.

\section{Natural translation}

The translation~$\nu$, defined below, transforms an NNF rule~(\ref{rule})
into a formula over the signature~$\sigma^\#$ that may include second-order
variables.  We define first how~$\nu$ applies to subexpressions of the rule.

\medskip\noindent 1. \emph{Applying~$\nu$ to terms.}
Let~$V_1,\dots,V_m$ be all variables that occur in~(\ref{rule}) at least once
\begin{itemize}
\item
  in the scope of an arithmetic function symbol, or
\item
  in the left-hand side of an equation containing the interval symbol.
\end{itemize}
Designate~$m$ distinct integer variables as $\nu(V_1),\dots,\nu(V_m)$.
If~$t$ is a term that does not contain~$..$ then $\nu(t)$ is the term over
the signature~$\sigma^\#$ obtained from~$t$ by
\begin{itemize}
\item substituting $\nu(V_i)$ for each variable~$V_i$, and
\item replacing each symbolic constant~$c$ in a subexpression of the
  form $c(t_1,\dots,t_k)$ by~$c\bs k$.
\end{itemize}

\medskip\noindent 2. \emph{Applying~$\nu$ to literals:}
\begin{itemize}
\item
$\nu(p({\bf t}))$ is $(p/k)(\nu({\bf t}))$,
\item $\nu(\no\ p({\bf t}))$ is $\neg (p/k)(\nu({\bf t}))$,
\item $\nu(\no\ \no\ p({\bf t}))$ is $\neg\neg (p/k)(\nu({\bf t}))$,
\end{itemize}
where~$k$ is the length of the tuple~$\bf t$.

\medskip
The transformation~$\nu$ applies to tuples of terms componentwise.

\medskip\noindent \emph{3. Applying~$\nu$ to comparisons: }
$\nu(t_1\prec t_2)$ is $\nu(t_1)\prec\nu(t_2)$, if~$t_2$ does not
contain the interval symbol;
$\nu(t = t_1\,..\,t_2)$ is $\nu(t_1)\leq\nu(t)\leq \nu(t_2)$.

\medskip
The transformation~$\nu$ applies to tuples of literals and comparisons
componentwise.

\medskip\noindent \emph{4. Applying~$\nu$ to aggregate atoms and to negated
aggregate atoms.}  Make the list of all occurrences of
aggregate elements in the body of~(\ref{rule}) and choose, for each
aggregate element $X_1,\dots,X_k:{\bf L}$ on the list, a
variable~$C$ \emph{associated} with it, so that
\begin{itemize}
\item if {\bf L} contains no global variables
  then~$C$ is an object variable of the sort \emph{set};
\item if the global variables occurring in {\bf L} are $Y_1,\dots,Y_l$
  ($l>0$) then~$C$ is a function variable of arity~$n$; the sort of its value
  is \emph{set}; the sort of its $i$-th argument is \emph{integer}
  if~$Y_i$ is one of the variables $V_1,\dots,V_m$, and \emph{general}
  otherwise.
\end{itemize}
The formula $\sd(C)$  (``the set definition for~$C$'') is defined as follows.
If $C$ is an object variable then $\sd(C)$ is
$$\forall V(V\in C\lrar\exists\, \nu(X_1)\cdots\nu(X_k)\,
(V=(\nu(X_1),\dots,\nu(X_k))
\land \nu({\bf L}))),$$
where~$V$ is a general variable.
(If $k=1$ then $(\nu(X_1),\dots,\nu(X_k))$ is understood as~$\nu(X_1)$.)
If~$C$ is a function variable then $\sd(C)$ is
$$\ba l
\forall V \nu(Y_1) \cdots \nu(Y_l)(V\in C(\nu(Y_1),\dots,\nu(Y_l))\\
\qquad\lrar
\exists\, \nu(X_1)\cdots\nu(X_k)\,(V=(\nu(X_1),\dots,\nu(X_k))
\land \nu({\bf L}))).
\ea$$

(Set definitions $\sd(C)$ are not used in the process of applying~$\nu$ to a
rule, but they will be referred to in the definition of completion.)

The result of applying~$\nu$ to an aggregate atom
containing an aggregate element $X_1,\dots,X_k:{\bf L}$ is the formula
obtained by replacing $\co\{X_1,\dots,X_k:{\bf L}\}$ with
$$\co(C(\nu(X_1),\dots,\nu(X_k))).$$
Applying~$\nu$ to a negated aggregate atom is performed in the same way,
with $\no$ replaced by~$\neg$.

\medskip\noindent
\emph{5.} The result of applying~$\nu$ to the empty string is~$\bot$.

\medskip\noindent
\emph{6.} The result of applying~$\nu$ to rule~(\ref{rule}) is the formula
$$
\forall {\bf Z}(\nu(B_1)\land\cdots\land\nu(B_n) \rar \nu(H)),
$$
where $\bf Z$ is obtained from the list of global
variables occurring in the rule
by replacing $V_1,\dots,V_m$ with $\nu(V_1),\dots,\nu(V_m)$.  The free
variables of this formula are the variables associated with aggregate
expressions.

The result of applying $\nu$ to a semi-normal rule is defined as the result of applying $\nu$ to the NNF rule obtained from it as described at the end of Section~\ref{sec:nnf}.

% \section{Examples}

  

\bibliographystyle{named}
\bibliography{bib}
\end{adjustwidth}
\end{document}

